\documentclass[12pt]{article}
\usepackage[utf8]{inputenc}
\usepackage[T1]{fontenc}
\usepackage{hyperref}
\title{Hasha Language Resources for K-12 Curriculum Development}
\author{Compiled by Claude Sonnet 5.5 (claude-sonnet-5-5) for Andr\'as Kornai}
\date{}
\begin{document}
\maketitle

\section{Introduction}
Hasha (ISO 639-3 ybj; Glottocode hash1238; also called Yashi) is a Plateau language of Nigeria. Wikipedia and Omniglot place it in Nasarawa State, mainly Wamba Local Government Area, while one Endangered Languages record gives Akwanga LGA in Plateau State. Estimates of speakers conflict: about 3,000 on Wikipedia and Omniglot, about 400 on Verbix. A 1999 survey by Roger Blench found it still spoken by young people. Few speakers are literate in it, and many also speak Hausa and/or English. Very little material exists.

\section{Primary Educational Resources}
No primer, textbook or school programme for Hasha turned up in the searches.

\section{Linguistic and Reference Materials}
\begin{itemize}
\item Wikipedia article on the language: \url{https://en.wikipedia.org/wiki/Hasha_language}
\item Glottolog references by Roger Blench: \url{https://glottolog.org/resource/reference/id/599473}, \url{https://glottolog.org/resource/reference/id/150919}, \url{https://glottolog.org/resource/reference/id/120179}. I did not check the titles of these entries.
\item Endangered Languages record for the language: \url{https://endangeredlanguages.com/lang/485}
\item Endangered Languages record ``Field trip to record the status of some little-known Nigerian languages'': \url{https://endangeredlanguages.com/node/84505}
\end{itemize}

\section{Religious and Community Resources}
\begin{itemize}
\item Scripture Earth page for ybj. I did not verify whether any scripture exists in Hasha: \url{https://www.scriptureearth.org/00eng.php?iso=ybj}
\end{itemize}

\section{Digital and Audio Resources}
\begin{itemize}
\item Omniglot page on Hasha: \url{https://omniglot.com/writing/hasha.htm}
\item Verbix page for Hasha: \url{https://docs.verbix.com/EndangeredLanguages/Hasha}
\end{itemize}

\section{Curriculum Development Context}
Documentation is thin and no orthography in community use was confirmed. A curriculum would need basic documentation and community consultation first, and the conflicting speaker counts should be resolved.

\end{document}